\section*{第 \ref{ch:g7:shadows-eclipses} 章——影子与日食}

\begin{solution}{exo:g7:shadows-eclipses:1}
\omterm{prop:g7:shadows-eclipses:umbra}{本影}：\omterm{def:g1:light-and-shadows:source}{光源}被完全遮住，\omterm{def:g1:light-and-shadows:source}{光源}任何部分的光都到达不了。
\omterm{prop:g7:shadows-eclipses:umbra}{半影}：\omterm{def:g1:light-and-shadows:source}{光源}被部分遮住，\omterm{def:g1:light-and-shadows:source}{光源}的一部分光仍可越过遮挡物的边缘
照过来。
\end{solution}

\begin{solution}{exo:g7:shadows-eclipses:2}
点\omterm{def:g1:light-and-shadows:source}{光源}没有可以被部分遮住的“部分”：每处位置要么看得见它，
要么看不见它——\omterm{def:g1:light-and-shadows:shadow}{影子}轮廓分明，没有\omterm{prop:g7:shadows-eclipses:umbra}{半影}。阴天的天空是一个很大的\omterm{def:g1:light-and-shadows:source}{光源}，
几乎每处都能看到它的大部分：
\omterm{def:g1:light-and-shadows:shadow}{影子}几乎全是\omterm{prop:g7:shadows-eclipses:umbra}{半影}，被冲得模糊不清。
\end{solution}

\begin{solution}{exo:g7:shadows-eclipses:3}
从\omterm{def:g1:light-and-shadows:source}{光源}的\emph{边缘}发出、擦过遮挡物边缘的\omterm{def:g6:light-propagation:ray}{光线}，
两条边缘\omterm{def:g6:light-propagation:ray}{光线}都照不到的区域即为\omterm{prop:g7:shadows-eclipses:umbra}{本影}。
\end{solution}

\begin{solution}{exo:g7:shadows-eclipses:4}
\omterm{def:g7:shadows-eclipses:solar}{日食}：太阳—\omterm{def:g2:moon-in-the-sky:moon}{月球}—地球，顺序如此——\omterm{def:g2:moon-in-the-sky:moon}{月球}必须是新月
（位于我们与太阳之间）。\omterm{def:g7:shadows-eclipses:lunar}{月食}：太阳—地球—\omterm{def:g2:moon-in-the-sky:moon}{月球}——\omterm{def:g2:moon-in-the-sky:moon}{月球}必须是\omterm{ex:g2:moon-in-the-sky:shapes}{满月}
（位于太阳的对面）。
\end{solution}

\begin{solution}{exo:g7:shadows-eclipses:5}
全食：只有处在\omterm{prop:g7:shadows-eclipses:umbra}{本影}那一小块快速移动的光斑之内的人才能看到。
偏食：范围很大的\omterm{prop:g7:shadows-eclipses:umbra}{半影}周围区域。完全看不到：
所有\omterm{def:g1:light-and-shadows:shadow}{影子}区域都照不到的人——即地球上的大部分地区。
\end{solution}

\begin{solution}{exo:g7:shadows-eclipses:6}
\omterm{def:g7:shadows-eclipses:lunar}{月食}是\omterm{def:g2:moon-in-the-sky:moon}{月球}本身变暗——这是整个夜半球的人都能同时看到的变化。
日全食则要求站在\omterm{def:g2:moon-in-the-sky:moon}{月球}细长\omterm{prop:g7:shadows-eclipses:umbra}{本影}的\emph{内部}，
即\omterm{prop:g7:shadows-eclipses:umbra}{本影}触及地面的地方，
那只是一个光斑，而非覆盖全天的事件。
\end{solution}

\begin{solution}{exo:g7:shadows-eclipses:7}
\omterm{def:g2:moon-in-the-sky:moon}{月球}\omterm{def:g6:sun-earth-moon:axisorbit}{轨道}相对地球—太阳所在平面倾斜约五度，
因此大多数新月和\omterm{ex:g2:moon-in-the-sky:shapes}{满月}都从精确的连线上方或下方经过，
\omterm{def:g1:light-and-shadows:shadow}{影子}便错过了。只有当月相与穿越平面的时刻恰好重合时，
三者才会排成一线。
\end{solution}

\begin{solution}{exo:g7:shadows-eclipses:8}
用针孔把太阳投影到纸上，以及树叶在地面形成的
月牙形光斑（或用漏勺）；直接观看时须戴合格、完好的\omterm{def:g7:shadows-eclipses:solar}{日食}眼镜。
裸眼只能在全食期间观看，
一旦出现第一缕重新露出的太阳，就须停止观看。
\end{solution}

\begin{solution}{exo:g7:shadows-eclipses:9}
太阳的宽度约为\omterm{def:g2:moon-in-the-sky:moon}{月球}的 $400$ 倍，距离约为其
$400$ 倍，因此两者的圆盘在天空中几乎完全吻合。
恰好吻合。这种吻合的结果是：日全食确实存在，但只是勉强——
时间很短，\omterm{prop:g7:shadows-eclipses:umbra}{本影}光斑也很小；当\omterm{def:g2:moon-in-the-sky:moon}{月球}运行到椭圆\omterm{def:g6:sun-earth-moon:axisorbit}{轨道}较远的一段时，
它的圆盘略小，无法完全遮住太阳的边缘，
于是留下一圈未被遮住的光环——即日环食。
\end{solution}

\begin{solution}{exo:g7:shadows-eclipses:10}
每个\omterm{def:g1:light-and-shadows:shadow}{影子}的\omterm{prop:g7:shadows-eclipses:umbra}{半影}边界都由太阳两侧的边缘决定，作图表明，
\omterm{prop:g7:shadows-eclipses:umbra}{半影}随遮挡物与屏幕之间的间距增大而加宽：
来自太阳相对两侧边缘的\omterm{def:g6:light-propagation:ray}{光线}，在越过遮挡物后走得越远，
彼此分开得越远。脚下的间距为零，\omterm{def:g1:light-and-shadows:shadow}{影子}边界清晰；
头部所对应的\omterm{def:g1:light-and-shadows:shadow}{影子}处在相距两米、已经扩散开的位置，
边界柔和。
\end{solution}

\begin{solution}{exo:g7:shadows-eclipses:11}
每道树叶缝隙都是天然的针孔相机，每块地面光斑都是
该针孔所成的太阳的像——平常的日子里，它是圆形的。
偏食期间，太阳本身呈月牙形，因此每台忠实的小相机都投影出月牙形，
而且它们全都指向同一方向，因为它们都在复制天空中
那一个月牙形。
\end{solution}

\begin{solution}{exo:g7:shadows-eclipses:12}
掠过地球边缘的阳光沿\omterm{def:g5:air-pressure-weather:pressure}{大气层}中尽可能长的路径穿行——
同时在整个地球周围形成与\omterm{def:g1:day-and-night:daynight}{日落}时相当的长路径。
\omterm{def:g2:air-around-us:air}{空气}会将穿过的大部分光漫射掉（正是这部分光使地球的天空呈现颜色），
而漫射最先夺走的是路径较短的那些色光，
因此，穿过这段长途掠过地球的旅程、最终抵达\omterm{def:g2:moon-in-the-sky:moon}{月球}的光，
正是\omterm{def:g1:day-and-night:daynight}{日落}时幸存下来的那一部分：红色。
被食的\omterm{def:g2:moon-in-the-sky:moon}{月球}所发出的光，是地球上每一处\omterm{def:g1:day-and-night:daynight}{日出}和\omterm{def:g1:day-and-night:daynight}{日落}的光
同时投射到它身上的总和。
\end{solution}

\begin{solution}{pb:g7:shadows-eclipses:1}
\textbf{1.} 条纹是\omterm{def:g2:moon-in-the-sky:moon}{月球}\omterm{prop:g7:shadows-eclipses:umbra}{本影}掠过地面时的路径，
两侧宽阔的区域是\omterm{prop:g7:shadows-eclipses:umbra}{半影}
在地面上的投影。
\textbf{2.} 偏食：小镇位于\omterm{prop:g7:shadows-eclipses:umbra}{半影}之中——
太阳被咬去一块，应借助合格的防护用品观看，
但不会出现黑暗。
\textbf{3.} 即使只剩百分之五的太阳，其亮度仍是\omterm{ex:g2:moon-in-the-sky:shapes}{满月}光的数千倍：
它足以让天空保持明亮，并使日冕不可见。
正午如夜的景象只属于\omterm{prop:g7:shadows-eclipses:umbra}{本影}，
不存在“几乎全食”这种情况。
\textbf{4.} \omterm{prop:g7:shadows-eclipses:umbra}{本影}光斑大致为圆形：位于中线上的营地在光斑
掠过时处于光斑的全部宽度之下，
而靠近条纹边缘的营地只会经过圆的一条弦——
经过的路径较短，全食的时间也就较短。
\textbf{5.} 合格的\omterm{def:g7:shadows-eclipses:solar}{日食}眼镜能挡住太阳的灼烧\omterm{def:g6:light-propagation:ray}{光线}；
太阳镜则不能——它只能减弱眩光，
却会让灼伤视网膜的那部分光透过，
使人能够舒适地长时间直视太阳，比不戴更糟糕。
\textbf{6.} 一端开一个小孔，另一端放置描图纸作屏幕
（或直接用地面上的白纸）：让\emph{背}对着太阳站立，
将小孔对准肩后方，观察投影出的圆盘——
切勿透过小孔直视太阳。
\textbf{7.} 不可取：水面的\omterm{def:g5:mirrors-reflection:reflection}{反射}是太阳的镜像，
其刺眼程度几乎与太阳本身相当——经\omterm{def:g5:mirrors-reflection:reflection}{反射}后的灼伤
仍然是灼伤。
\textbf{8.} 纸上有几十个月牙形的像：每个小孔都是一台针孔相机，
每台都投射出被咬去一块的太阳——
漏勺里装满了\omterm{def:g7:shadows-eclipses:solar}{日食}。
\textbf{9.} \omterm{def:g2:moon-in-the-sky:moon}{月球}在强光的映衬下原本不可见，直到它的圆盘
边缘开始遮住太阳的边缘——营地此时已进入
\omterm{prop:g7:shadows-eclipses:umbra}{半影}之中。
\textbf{10.} 可见的太阳正在收缩成一道细细的弧，
——\omterm{def:g1:light-and-shadows:source}{光源}变得越来越像一个点。\omterm{def:g1:light-and-shadows:source}{光源}越小，\omterm{prop:g7:shadows-eclipses:umbra}{半影}边界越窄：
每个\omterm{def:g1:light-and-shadows:shadow}{影子}的柔边都在收紧，
周围的世界变得奇异而分明。
\textbf{11.} \omterm{def:g2:moon-in-the-sky:moon}{月球}的\omterm{prop:g7:shadows-eclipses:umbra}{本影}本身——全班同学此时站在影锥的顶端，
随着\omterm{def:g1:light-and-shadows:shadow}{影子}掠过营地。黑色的圆盘周围悬挂着日冕，
即太阳珍珠般的外层大气，此刻才能看见；
第一道重新露出的太阳边缘恢复了正常的白昼光照——
请戴上眼镜。
\textbf{12.} 例如：“\omterm{def:g7:shadows-eclipses:solar}{日食}是一种\omterm{def:g1:light-and-shadows:shadow}{影子}，它的\omterm{def:g1:light-and-shadows:source}{光源}、遮挡物和屏幕
都是天体：\omterm{def:g2:moon-in-the-sky:moon}{月球}的\omterm{prop:g7:shadows-eclipses:umbra}{本影}掠过我们脚下的大地，
或是\omterm{def:g2:moon-in-the-sky:moon}{月球}穿过我们的\omterm{def:g1:light-and-shadows:shadow}{影子}——这是一种
只需两条\omterm{def:g6:light-propagation:ray}{光线}和一把直尺便能画出的几何，
在自然界中以最宏大的尺度上演。”
\end{solution}
